\section{Fields Across Boundaries}
\label{sec:ch06-fields-across-boundaries}
\index{federation!field-level directives|(}%

Four directives describe what happens to a field when the type it hangs on lives
in more than one subgraph. None of them make sense on a value type, because a
value type belongs to exactly one subgraph and has no second place for a field
to come from.

\subsection*{\code{@shareable}: a field more than one subgraph can resolve}

\index{federation!shareable@\code{@shareable}}%
Without \code{@shareable}, a given field of a given object type may appear in
exactly one subgraph~\autocite{apollo2026compositionrules}. Put \code{Product.title} in Catalog and again in Reviews,
without the directive, and composition fails. The rule exists because Federation
v1 did not have it, and the result was a schema that looked correct and
sometimes produced answers from the wrong service. \code{@shareable} makes the
intention explicit: this field can be resolved by more than one subgraph, and
the router may pick either.

\begin{graphqlsdl}
# In Catalog
type Product @key(fields: "id") {
  id: ID!
  title: String! @shareable
  price: Money!
}

# In Inventory
type Product @key(fields: "id") {
  id: ID!
  title: String! @shareable
  stock: Int!
}
\end{graphqlsdl}

Here \code{title} appears in both, and the router can fetch it from whichever it
is already talking to. The practical use is reducing round trips: if a query
selects \code{price} and \code{title} and \code{stock}, the router might fetch
\code{title} from Inventory alongside \code{stock} rather than making a third
call to Catalog for one field. But that is a query-planning decision the router
makes, not a guarantee you write into the schema, and
chapter~\ref{ch:how-a-federated-query-actually-runs} shows the plan that
resulted.

The cost is that both subgraphs now have to keep the field correct. If Catalog
updates \code{title} and Inventory does not, the router serves stale data with
no way to know it is stale. That is not a bug in the router. It is the
definition of a shareable field: two sources of truth that agree because you
keep them agreeing, not because the system enforces it.

\subsection*{\code{@external}: this field comes from somewhere else}

\index{federation!external@\code{@external}}%
\code{@external} marks a field as defined elsewhere in the supergraph. It
appears in a subgraph alongside \code{@requires} or \code{@provides}, never
alone. The field still appears on the type, and the subgraph still knows its
type, but the router does not ask this subgraph to resolve it.

\begin{graphqlsdl}
# In the Shipping subgraph
type Product @key(fields: "id") {
  id: ID!
  weight: Float! @external
  shippingCost: Money! @requires(fields: "weight")
}
\end{graphqlsdl}

Shipping does not own \code{weight}. Catalog does. But \code{shippingCost}
depends on it, so Shipping declares the field as \code{@external} so the router
knows to fetch it from Catalog first and hand it along. The word
\enquote{external} is from the subgraph's point of view: this field is part of
the type's schema, but it is resolved by somebody else.

\subsection*{\code{@requires}: I need these to compute mine}

\index{federation!requires@\code{@requires}}%
\code{@requires} names the external fields a resolver cannot function without.
The example above says \code{shippingCost} needs \code{weight}. When a query
selects \code{shippingCost}, the router fetches \code{weight} from Catalog first
and includes it in the representation sent to Shipping. The resolver receives
the \code{weight} and computes the cost, and it never had to ask for it.

The \code{fields} argument is a GraphQL selection set, not just a flat list. It
can reach into nested types:

\begin{graphqlsdl}
type Product @key(fields: "id") {
  id: ID!
  dimensions: Dimensions @external
  volumetricWeight: Float! @requires(fields: "dimensions { length width height }")
}
\end{graphqlsdl}

\index{federation!requires cost@\code{@requires} cost}%
The router fetches the whole selection before calling the resolver. That is the
performance guarantee and the cost. A resolver that needs one small field costs
one extra subgraph fetch for that field, and a resolver that needs a deep
traversal costs the same traversal. Chapter~\ref{ch:entity-resolution-done-right}
measures the difference.

\subsection*{\code{@provides}: I can give you extra fields from this entity}

\index{federation!provides@\code{@provides}}%
\code{@provides} is the inverse of \code{@requires} and the more surprising of
the two. It says that when this field returns an entity, the resolver can also
supply the named fields of that entity, saving the router a follow-up fetch.

\begin{graphqlsdl}
# In the Reviews subgraph
type Review @key(fields: "id") {
  id: ID!
  body: String!
  author: Customer! @provides(fields: "displayName")
}

type Customer @key(fields: "id") {
  id: ID!
  displayName: String! @external
}
\end{graphqlsdl}

Reviews returns a \code{Customer} as part of a \code{Review}'s author. The
router would normally fetch \code{displayName} from Accounts separately.
\code{@provides(fields: "displayName")} tells the router that the resolver
already has it, so the router can skip the call. The subgraph is making a
promise it has to keep: if it says it provides \code{displayName}, the resolver
had better include it in the returned \code{Customer}, or the client gets null
and the router did not ask anybody else.

\index{federation!provides vs shareable@\code{@provides} vs \code{@shareable}}%
This is not the same as \code{@shareable}. \code{@shareable} says a field can
live in two subgraphs. \code{@provides} says a resolver for one field can supply
another field of the entity it returns, and that other field is still owned by
one subgraph. The router receives it from Reviews and does not fetch it from
Accounts, but Accounts is still where it is defined.

The directive is an optimisation and it is the one that most often goes wrong.
Getting it right means the resolver returns more than it was asked for, which
GraphQL deliberately does not do anywhere else. Chapter~\ref{ch:entity-resolution-done-right}
has the measured numbers and the failure modes.

\index{federation!field-level directives|)}%
